%% Generated by verification/make_figures.py from regular-graph-harmonic-quotient-examples.tex.
%% Figure fig:harmonic-quotient_d5reg10d+St6 (replaces figures/harmonic-quotient_d5reg10d+St6-G2.tex).
%% Gamma_2 = d5reg10d: 10 vertices, 25 edges (edge list in the list item).
%% phi from the table at line 9834. Vertex x is filled with color hq<phi(x)>.
%% Layout: hand-tuned.
\begin{tikzpicture}[hqgraph]
\node[hqv, fill=hq4] (v0) at (-1.770,-1.416) {0};
\node[hqv, fill=hq1] (v1) at (-1.770,1.416) {1};
\node[hqv, fill=hq1] (v2) at (-0.885,1.416) {2};
\node[hqv, fill=hq0] (v3) at (-0.885,-1.416) {3};
\node[hqv, fill=hq3] (v4) at (0.000,-1.416) {4};
\node[hqv, fill=hq1] (v5) at (0.000,1.416) {5};
\node[hqv, fill=hq1] (v6) at (0.885,1.416) {6};
\node[hqv, fill=hq5] (v7) at (0.885,-1.416) {7};
\node[hqv, fill=hq1] (v8) at (1.770,1.416) {8};
\node[hqv, fill=hq2] (v9) at (1.770,-1.416) {9};
\draw[hqe] (v0) -- (v1);
\draw[hqe] (v0) -- (v2);
\draw[hqe] (v0) -- (v5);
\draw[hqe] (v0) -- (v6);
\draw[hqe] (v0) -- (v8);
\draw[hqe] (v1) -- (v3);
\draw[hqe] (v1) -- (v4);
\draw[hqe] (v1) -- (v7);
\draw[hqe] (v1) -- (v9);
\draw[hqe] (v2) -- (v3);
\draw[hqe] (v2) -- (v4);
\draw[hqe] (v2) -- (v7);
\draw[hqe] (v2) -- (v9);
\draw[hqe] (v3) -- (v5);
\draw[hqe] (v3) -- (v6);
\draw[hqe] (v3) -- (v8);
\draw[hqe] (v4) -- (v5);
\draw[hqe] (v4) -- (v6);
\draw[hqe] (v4) -- (v8);
\draw[hqe] (v5) -- (v7);
\draw[hqe] (v5) -- (v9);
\draw[hqe] (v6) -- (v7);
\draw[hqe] (v6) -- (v9);
\draw[hqe] (v7) -- (v8);
\draw[hqe] (v8) -- (v9);
\end{tikzpicture}%
